% !TeX root = ../教材题库_全部条目_转写工作稿.tex
% Source: XJ-B1-C05-EX-P199-Y0577
% 第5章 / 5.5 三角函数模型的简单应用 / 教材 p.199 / PDF 第206页
% OCR未校对：公式与文字均可能有误；未完成数学TeX转写；禁止直接发给学生
% 原题及相邻原解截图（核对用）：../题目原图/XJ-B1-C05-EX-P199-Y0577-01.webp
% 本文件只在转写完成后作为题干使用；OCR 条目尚不是可靠题干。
\begin{bankitem}{5.5 例题 例题1}{第5章 / 5.5 三角函数模型的简单应用 / 教材 p.199 / PDF 第206页}{XJ-B1-C05-EX-P199-Y0577}{OCR未校对：公式与文字均可能有误；未完成数学TeX转写；禁止直接发给学生}
\begin{ocrdraft}
[教材p.199 / PDF第206页]
\par
00图5，5一1为小球在做单摆运动（可近似看作简谐振动）时，离开平衡位置
\par
时的位移）随时间）变化的函数图象，已知该图象满足y=Asin()x十f)(xe
\par
[0，十。。），仞>0，0<<爿的形式．试根据函数图象求出这个单摆运动的函数解
\par
析式．
\par
2
\par
0
\par
5訌
\par
12
\par
12
\par
一2
\par
图5．5
\par
11丌5丌
\par
解由图象知，周期T=2（一一
\par
1212
\par
因为点（老
\par
0）在函数图象上，且函数图象满足3，一Asin（-十@的形式，
\par
5丌
\par
5丌
\par
所以Asin2\ensuremath{\times}軛十=0，即sin一十=0．
\par
5丌5丌
\par
4丌
\par
又已知0<<《，则<十<
\par
5丌
\par
从而十一丌，即一否．
\par
本例还可以利用关键
\par
又点（0，1)在函数图象上，所以Asin
\par
=1，得A
\par
=2．
\par
点结合待定系数法来解
\par
决，自己动手试一试，
\par
故所求函数的解析式为Y
\par
=2sin2x十
\par
6
\end{ocrdraft}
\end{bankitem}
